\par\Needspace{14\baselineskip}

\label{sec:analysis-multivariable-differentiation}
\label{sec:multivariable-differentiation}

\section{多元函数与场}
\label{sec:multivariable-functions-and-fields}

微分和积分都要先说清研究的是什么函数。本部分最常见的情形是：空间中每个位置都带着一个可以测量的量。平板上每一点有温度，河水中每一点有流速，金属块中每一点有密度。这样的量随位置变化，输入是位置的坐标，输出是该位置上读出的量。

第~\ref{sec:multivariable-joint-limit-section}节已经从极限的角度定义了多元标量函数与向量值映射。本节换一个角度，把它们看作空间位置上的属性：先按输出是一个数还是一个向量，区分标量场与向量场；再说明本部分的定义域、输出空间以及极限与连续性怎样理解。这些约定确定了后面各章微分与累计的对象。

\subsection{标量场}
\label{sec:multivariable-scalar-fields}

温度计放在平板的不同位置，读数不同；把每个位置与它的读数配对，就得到一个函数。密度、压强、电势也都在每个位置给出一个实数。这类函数在每个位置只给出大小，不带方向。
\begin{definition}[标量场]
\label{def:multivariable-scalar-field}
设$\Omega\subseteq\mathbb R^d$。映射$f:\Omega\to\mathbb R$称为$\Omega$上的\term{标量场}（scalar field）；它把每个位置$\bm x\in\Omega$映为一个实数$f(\bm x)$。
\end{definition}
作为映射，标量场就是定义~\ref{def:multivariable-scalar-function}中的多元标量函数。“场”这个名称只强调输入是空间位置、输出是该位置的属性。后面称它为多元函数还是标量场，取决于讨论的是一般的输入输出关系，还是位置上的物理量，数学内容相同。

以$f(x,y)=x^2+y^2$为例，它在每个位置给出到原点距离的平方。它有两种画法。第一种画法把函数值当作高度，得到三维空间中的曲面$z=x^2+y^2$，形状像一只碗。第二种画法留在输入平面里，画出等值曲线$\{(x,y):x^2+y^2=c\}$。

对这个函数，$c>0$时等值曲线是半径为$\sqrt c$的圆，$c=0$时只剩原点，$c<0$时是空集。取$c=1,2,3,4$，半径依次为$1,\sqrt2\approx1.41,\sqrt3\approx1.73,2$，相邻两圈的间距越来越小。函数值每增加$1$所需走的距离越来越短，说明离原点越远，场值沿径向增长得越快。

这一观察还只是定性的。度量“沿某方向增长多快”需要导数：第~\ref{sec:multivariable-gradient-inner-product}节把它写成梯度，并说明梯度与等值曲线正交。

\subsection{向量场}
\label{sec:multivariable-vector-fields}

河水在每一点有流速，物体在重力场或电场中的每一点受力。这些量既有大小又有方向，并且与位置同处一个空间：流速可以画成从该点出发的箭头。
\begin{definition}[向量场]
\label{def:multivariable-vector-field}
设$\Omega\subseteq\mathbb R^d$。映射$\bm F:\Omega\to\mathbb R^d$称为$\Omega$上的\term{向量场}（vector field）。按本书约定，$\bm F(\bm x)$写成$d\times1$列矩阵$(F_1(\bm x),\ldots,F_d(\bm x))^\top$，每个分量$F_i:\Omega\to\mathbb R$都是标量场。
\end{definition}
向量场的输入是位置，输出是放在该位置的一个向量。例如$\bm F(x,y)=(x,y)^\top$的箭头从原点向外，离原点越远越长；$\bm F(x,y)=(-y,x)^\top$的箭头与以原点为圆心的圆相切，按逆时针方向。第~\ref{chap:path-integrals}章用这两个场比较沿路径做的功。

向量场是向量值函数的特例。一般向量值函数$F:\Omega\subseteq\mathbb R^d\to\mathbb R^m$的输出维数$m$可以与输入维数不同：参数曲线把一个参数映为空间中的点，神经网络的一层把输入向量映为另一维数的向量。向量场要求$m=d$，输出向量才能画成输入空间中的箭头，也才能与路径的切向量或曲面的法向量作内积。

输出还可以是矩阵。若每个位置给出一个$p\times q$矩阵，例如介质中每一点的应力矩阵，就得到矩阵值的场；它的每个元素仍是标量场。第~\ref{chap:matrix-differentiation}章研究的则是另一种情形：输入本身是矩阵，例如$\bm X\mapsto\bm X^{-1}$。判断一个映射属于哪一类时，要分别核对输入空间和输出空间，不能只看表达式里是否出现矩阵。

\subsection{多元函数的极限与连续}
\label{sec:multivariable-field-limits-continuity}

标量场和向量场的极限与连续性已在第~\ref{chap:normed-inner-product-spaces}章建立，这里只回顾后面要用的结论，不重复证明。

极限用输入、输出空间的范数定义（定义~\ref{def:multivariable-limit}）。$\lim_{\bm x\to\bm a}F(\bm x)=\bm L$的意思是：只要输入进入以$\bm a$为心的某个去心小球，输出与$\bm L$的距离就小于给定误差，不论输入沿什么路线进入。连续性（定义~\ref{def:multivariable-scalar-continuity}）要求这个极限等于$F(\bm a)$。

对有限维的向量输出，这些判断可以逐分量进行：向量场或矩阵值的场在一点有极限或连续，当且仅当每个输出分量如此（定理~\ref{thm:multivariable-coordinatewise-convergence}）。逐分量分解的是输出。输入的多个坐标不能这样拆开，因为输入可以同时改变全部坐标、沿弯曲的路线趋近。

因此，沿路径检验只能用来否定极限。沿两条路径得到不同的极限值，联合极限就不存在：示例~\ref{ex:multivariable-path-limits}中的$xy/(x^2+y^2)$沿$x$轴趋于$0$，沿$y=x$趋于$1/2$。反过来，只检查若干条路径不能证明极限存在：示例~\ref{ex:multivariable-parabola-path-limit}中的$x^2y/(x^4+y^2)$沿每条直线都趋于$0$，沿抛物线$y=x^2$却恒为$1/2$。要证明极限存在，需要一个只依赖输入到基点距离的统一误差估计。

这个区别在微分中还会出现。方向导数是沿固定直线取的极限，全微分则要求误差对全部趋近方式统一受控；第~\ref{sec:multivariable-total-differentiation}节的反例与上面的抛物线例子结构相同。

最后约定定义域。讨论微分时，基点取在开集内部，从基点出发向任何方向走足够小的一步都不离开定义域；讨论积分时，累计范围通常是有界闭区域，这时只要求被积函数在区域上连续或满足相应的可积条件。
